\documentclass[a4paper,twoside]{article}

\usepackage{morewrites}

\usepackage[svgnames,dvipsnames,x11names]{xcolor}
\usepackage{graphicx}
\usepackage[english]{babel}
\usepackage[colorlinks=true, linkcolor=NavyBlue, urlcolor=NavyBlue, citecolor=NavyBlue]{hyperref}
\usepackage[a4paper]{geometry}
\usepackage{ts-fonts}
\usepackage{progressbar}
\usepackage{amsmath}
\usepackage{float}
\usepackage{seqsplit}
\usepackage{ts-code}

\usepackage{lastpage}
\usepackage{refcount}
\makeatletter
\def\ps@plain{
  \let\@oddhead\@empty
  \let\@evenhead\@empty
  \def\@oddfoot{
    \hfil
    \ifnum\getpagerefnumber{LastPage}>1\relax
      \thepage
    \fi
    \hfil}
  \let\@evenfoot\@oddfoot
}
\makeatother

\title{Progress Bars}
\author{}\date{}

\begin{document}

\maketitle

\thispagestyle{plain}
The \texttt{texsmith.progressbar} extension renders Markdown shorthand into \LaTeX{} progress bars:

\begin{code}{text}{}{}
\begin{Verbatim}[commandchars=\\\{\},breaklines, breakanywhere, commandchars=\\\{\}]
[=25\PYZpc{} \PYZdq{}Research\PYZdq{}]
[=50\PYZpc{} \PYZdq{}Implementation\PYZdq{}]
[=75\PYZpc{} \PYZdq{}Review\PYZdq{}]\PYZob{}: .candystripe\PYZcb{}
[=100\PYZpc{} \PYZdq{}Launch\PYZdq{}]\PYZob{}: .thin\PYZcb{}
\end{Verbatim}
\end{code}
{\progressbar[width=9cm,heighta=12pt,roundnessr=0.1,borderwidth=1pt,linecolor=black,filledcolor=black!60,emptycolor=black!10]{0.25} Research}\par

{\progressbar[width=9cm,heighta=12pt,roundnessr=0.1,borderwidth=1pt,linecolor=black,filledcolor=black!60,emptycolor=black!10]{0.5} Implementation}\par

{\progressbar[width=9cm,heighta=12pt,roundnessr=0.1,borderwidth=1pt,linecolor=black,filledcolor=black!60,emptycolor=black!10]{0.75} Review}\par

{\progressbar[width=9cm,heighta=6pt,roundnessr=0.1,borderwidth=1pt,linecolor=black,filledcolor=black!60,emptycolor=black!10]{1} Launch}\par

\end{document}
